\documentclass[11pt]{article}

\usepackage[a4paper,margin=27mm]{geometry}
\usepackage{amsmath,amssymb,amsthm,mathtools}
\usepackage{array,booktabs}
\usepackage{microtype}
\usepackage[hidelinks]{hyperref}

\newtheorem{theorem}{Theorem}
\newtheorem{proposition}[theorem]{Proposition}
\newtheorem{corollary}[theorem]{Corollary}
\newtheorem{lemma}[theorem]{Lemma}
\theoremstyle{definition}
\newtheorem{conjecture}[theorem]{Conjecture}
\newtheorem{remark}[theorem]{Remark}

\newcommand{\R}{\mathbb{R}}
\newcommand{\disc}{\Delta}
\newcommand{\eps}{\varepsilon}
\newcommand{\W}{W}

\title{Sharp Machine Thresholds, $1+\ln 2$ Amplification, BFS-Box Trapping ($\kappa_m\le m-1$),\\
and a Completeness Barrier for Machine-Dependent Weighted Chairman Assignment}
\author{Angus Muffatti}
\date{29 November 2026 --- research draft}

\begin{document}
\maketitle

\begin{abstract}
For a fractional assignment $x\in[0,1]^{m\times n}$ ($\sum_{i=1}^m x_{ij}=1$) with strictly positive machine-dependent weights $d\in(0,D]^{m\times n}$, let $\kappa_m$ denote the supremum over $m$-row instances of $\min_y \max_{i,t} |\sum_{j=1}^t d_{ij}(x_{ij}-y_{ij})|/D$ over all integral assignments $y\in\{0,1\}^{m\times n}$ ($\sum_i y_{ij}=1$, off-support allowed), and let $\kappa=\sup_{m\ge1}\kappa_m$.
Liu and Reis (ITCS 2026) proved the common-weight bound $\kappa(1)=1$ and conjectured $\kappa\le 1$ for machine-dependent weights; prior work in this programme established $\kappa_2=1$, refuted $\kappa_{11}\le 1$, and proved $\kappa\ge 1+\W(1/e)\approx 1.278465$ via single-level detector chains while documenting a composition barrier across blocks.

We settle the minimal machine threshold and substantially advance both lower and obstruction bounds:
\begin{enumerate}
\item \textbf{Exact phase transition at $m^*=3$:} using a zero-pollution initial split and dirty-helper ping-pong amplification, we prove $\kappa_3\ge c_3^*\approx 1.101721$ (the unique root of $2c-1=e^{-(c-1)}+e^{1-2c}$), alongside an explicit cleaned-helper $3\times 33$ counterexample ($\kappa_3>1001/1000$) certified in exact rational arithmetic, by independent exhaustive BFS over all $3^{33}$ schedules, and by HiGHS MILP. Combined with $\kappa_2=1$, the Liu--Reis bound holds on two machines and fails already on three. For every even $m\ge 4$, a zero-pollution ladder proves $\kappa_m\ge c_m^*$ solving $e^{-(c-1)}-(c-1)=2/m$ (giving $\kappa_4\ge 1.26625$ and $\kappa_8\ge 1.41220>1+1/e$).
\item \textbf{Breaking the composition barrier to $\kappa\ge 1+\ln 2\approx 1.693147$:} amplifying height \emph{inside} each fresh block eliminates entry pollution and accumulator venting, yielding three escalating support-2 families: fresh cascades ($1+1/e$), two-row ping-pong ($1+\W(1)\approx 1.567143$), and recursive binary-tree splitting plus ping-pong ($1+\ln 2\approx 1.693147$).
\item \textbf{Barrier for fresh-partner chains and abstract disjoint splits at $1+\ln 2$:} we prove universal optimality of the drain--gain law over arbitrary (possibly non-monotone) weights and splits, and construct a $1$-Lipschitz row-excess potential showing that no composition of forced fresh-partner chains ($\beta\le 0$) and abstract $q$-way disjoint-subtree splits ($q\ge 2$) reaches the single-row ping-pong threshold $h^*$ at any $c\ge 1+\ln 2$.
\item \textbf{Tightened universal upper bound $\kappa_m\le m-1$ via BFS-Box Trapping:} improving on the $0/1$-simplex baseline $\kappa_m\le m$ of B\'ar\'any--Grinberg / Lenstra--Shmoys--Tardos / Beck--Fiala, we prove via an oriented-pivot basic-feasible-solution trapping invariant that every $m$-row instance admits a support-preserving integral assignment with $|\disc_i(t)|\le (m-1)D$ (computable in $O(n\cdot\mathrm{poly}(m))$ exact arithmetic). This unifies $\kappa_2=1$ without 2D geometry, squeezes $1.10172\le \kappa_3\le 2$ and $1.26625\le \kappa_4\le 3$, restores the origin-enclosing convex-hull property $0\in\mathrm{conv}(R_t)$ at budget $c=m-1$, and is sharp for BFS-box trapping. We also resolve the first half of Open Problem~3 of the companion detector-chains draft via an explicit nonnegative dual identity.
\end{enumerate}
\end{abstract}

\section{Introduction}

Let $x\in[0,1]^{m\times n}$ satisfy $\sum_{i=1}^m x_{ij}=1$ for each column $j\in\{1,\dots,n\}$, and let $d\in(0,D]^{m\times n}$ be a matrix of strictly positive machine-dependent weights (normalized to $D=\max_{i,j}d_{ij}=1$ throughout). For an integral assignment $y\in\{0,1\}^{m\times n}$ with $\sum_{i=1}^m y_{ij}=1$ (not restricted to the support of $x$), define the row-prefix discrepancy
\[
 \disc_i(t)=\sum_{j=1}^t d_{ij}(x_{ij}-y_{ij}),\qquad
 \kappa_m=\sup_{m\text{-row instances}}\min_y\max_{i,t}\frac{|\disc_i(t)|}{D},\qquad
 \kappa=\sup_{m\ge1}\kappa_m.
\]
In the unweighted setting ($d_{ij}\equiv 1$), Tijdeman \cite{Tijdeman1973,Tijdeman1980} and Meijer \cite{Meijer1973} proved the exact optimal constant $1-\frac{1}{2m-2}$, hence $\kappa_{\text{unw}}=1$. Liu and Reis \cite{LiuReis2026} extended the guarantee $1$ to common column weights ($d_{ij}=d_j$, confirming a special case of the Morell--Skutella unsplittable-flow conjecture \cite{MorellSkutella2022}) and conjectured that $\kappa\le 1$ persists for arbitrary machine-dependent weights (Conjectures~19 and~21 in ITCS 2026; Conjectures~3 and~5 in arXiv:2511.18546).

Prior papers in this research programme established three landmarks:
\begin{itemize}
\item \texttt{chairman-counterexample/} refuted $\kappa_{11}\le 1$ via an explicit $11\times 15$ instance and proved $\kappa\ge 4-2\sqrt{2}\approx 1.171573$;
\item \texttt{chairman-small-rows/} proved $\kappa_2=1$ (via a parallelogram trapping invariant adapting Lund--Magazinov) and showed that the online problem is unbounded already for $m=2$;
\item \texttt{detector-chains/} generalized the detector gadget to depth-$k$ chains, proving $\kappa\ge 1+\W(1/e)\approx 1.278465$, while documenting a sharp composition barrier (entry pollution, accumulator venting, and disjunctive exits) that defeated every attempt to compose stronger local block games across blocks.
\end{itemize}

Note that every upper-sided, $\eta$-robust lower bound proved here transfers monotonically from $m$ to $m+1$ machines by padding with a dummy row carrying $x_{m+1,j}=0$ and $d_{m+1,j}=\eta\to 0^+$ (creating only an $O(n\eta)$ sink accommodated by every lower-bound proof below). Table~\ref{tab:summary} summarizes the main results of this paper (incorporating \texttt{frontier\_next/}, \texttt{frontier\_v2/}, and \texttt{frontier\_v3/} relative to the earlier \texttt{chairman-counterexample/}, \texttt{chairman-small-rows/}, and \texttt{detector-chains/} drafts).

\begin{table}[t]
\centering
\small
\begin{tabular}{@{}lll@{}}
\toprule
Quantity / Question & Prior Best & This Work \\
\midrule
Minimal $m^*$ with $\kappa_{m^*}>1$ & $m^*\in[3,11]$ & \textbf{$m^*=3$} (exact; $\kappa_2=1$, $\kappa_3\ge 1.101721$) \\
Smallest counterexample $(m,n)$ & $(11,15)$ & \textbf{$(3,33)$} (exact \texttt{Fraction} + MILP infeasible at $1$) \\
Small-$m$ lower bounds & $\kappa_2=1$, $\kappa_{11}>1$ & $\kappa_3\ge 1.10172$, $\kappa_4\ge 1.26625$, $\kappa_6\ge 1.36256$, $\kappa_8\ge 1.41220$ \\
Global lower bound on $\kappa$ & $1+\W(1/e)\approx 1.278465$ & \textbf{$1+\ln 2\approx 1.693147$} (via $1+1/e \to 1+\W(1) \to 1+\ln 2$) \\
Disjoint-tree barrier & Empirical ($c=1.35$–$1.60$) & \textbf{Proved fresh-partner / disjoint barrier at $1+\ln 2$} (Thms.~\ref{thm:nonmonotone}--\ref{thm:barrier}) \\
Upper bound on $\kappa_m$ & $2m$ ($\pm 1$ B\'ar\'any--Grinberg) & \textbf{$\kappa_m\le m-1$} (Thm.~\ref{thm:km1}: $\kappa_2\le 1$, $\kappa_3\le 2$, $\kappa_4\le 3$) \\
Detector-chains OP3 (part 1) & 2000-trial LP sweep & \textbf{Proved unconditionally} via exact dual (Thm.~\ref{thm:splitopt}) \\
\bottomrule
\end{tabular}
\caption{Summary of results on machine-dependent weighted chairman assignment.}
\label{tab:summary}
\end{table}

\section{The Drain--Gain Primitive and Internal Amplification to $\kappa\ge 1+\ln 2$}

Every lower-bound construction in this paper builds on a single elementary chain primitive. Instead of banking height \emph{across} blocks (where entry pollution and accumulator venting intervene), we amplify height \emph{inside each fresh block} before coupling to the global accumulator $B$. Every row inside a fresh block starts at height $\ge -n\eta$ (where $n$ is the total column count of the instance and $\eta>0$ is the off-support weight), and any off-support assignment on a column $C_j$ still denies both support rows their fractional shares while subtracting only $\eta$ from the recipient.

\subsection{Lemma A: Single-Level Drain--Gain Identity}

Fix a target budget $c>1$. Suppose an active row $X$ starts at height $\disc_X\ge\alpha$ with $0<c-\alpha<1$, and a passive row $Y$ starts at height $\disc_Y\ge\beta$. Append $k-1$ geometric chain columns $C_2,\dots,C_k$ supported on $\{X,Y\}$ with weights $d_j=r^{k-j}$ on $X$ ($d_k=1$, $r\in(0,1)$ chosen so $d_2=r^{k-2}\ge c-\alpha$), weight $1$ on $Y$, off-support weight $\eta$ elsewhere, and fractional splits
\[
 x_2=\frac{c-\alpha}{d_2}\in(0,1],\qquad x_j=r\in(0,1)\qquad (j=3,\dots,k).
\]

\begin{lemma}[Single-Level Drain--Gain]\label{lem:draingain}
Assume every prefix discrepancy stays within $[-c,c]$. If any column $C_j$ ($2\le j\le k$) is not assigned to $X$ (whether assigned to $Y$ or off-support), then at the first such column $C_j$, row $X$ reaches height at least $c$ (up to accumulated prior $\eta$-sinks). Otherwise all $k-1$ columns are assigned to $X$, and immediately after $C_k$:
\[
 \disc_X\ge \alpha-\sum_{j=2}^k d_j(1-x_j) = c-1,\qquad
 \disc_Y\ge \beta + 1 - x_2 + (k-2)(1-r).
\]
As $k\to\infty$ with $d_2\downarrow c-\alpha$ (so $x_2\uparrow 1$ and $(k-2)(1-r)\to-\ln(c-\alpha)$),
\[
 \disc_X \longrightarrow c-1,\qquad \disc_Y \longrightarrow \beta - \ln(c-\alpha).
\]
\end{lemma}
\begin{proof}
The drain sum telescopes identically for every $k\ge 3$:
\[
 \sum_{j=2}^k d_j(1-x_j)
 = d_2\left(1-\frac{c-\alpha}{d_2}\right)+\sum_{j=3}^k r^{k-j}(1-r)
 = (d_2 - c + \alpha) + (1 - d_2)
 = 1 - (c-\alpha).
\]
At the first column $C_j$ not assigned to $X$: for $j=2$, $\disc_X\ge\alpha+d_2 x_2=c$; for $j\ge 3$, the drain strictly before $C_j$ is $d_{j-1}-(c-\alpha)$ and $d_j x_j=d_{j-1}$, giving $\disc_X\ge\alpha-(d_{j-1}-c+\alpha)+d_{j-1}=c$. Sending $C_j$ off-support still denies $X$ and preserves these inequalities.
\end{proof}

\subsection{Fresh Cascades, Two-Row Ping-Pong, and Recursive Binary-Tree Splitting}

\begin{theorem}[$L$-Level Fresh Cascade]\label{thm:cascade}
For every $\eps>0$, there exists a finite support-2, strictly positive instance with $D=1$ forcing prefix discrepancy at least $1+1/e-\eps$. Hence $\kappa\ge 1+1/e\approx 1.367879$.
\end{theorem}
\begin{proof}
Fix $c\in(1,1+1/e)$. In each block $b=1,\dots,N$, introduce fresh rows $A_b,H_{b,1},\dots,H_{b,L}$ and place a detector $J_b$ on $\{B,A_b\}$ with weights $(1,1)$ and splits $(p,1-p)$, followed by $L$ levels $A_b\to H_{b,1}\to\cdots\to H_{b,L}$ designed via Lemma~\ref{lem:draingain} for the ideal heights $\alpha_0=1$, $\alpha_{r+1}=-\ln(c-\alpha_r)$. Since $\alpha+e^{-\alpha}\ge 1+1/e>c$ for all $\alpha\ge 1$, the increment $-\ln(c-\alpha)-\alpha$ is strictly positive and bounded below on $[1,c)$, so $\alpha_L>c$ for some finite $L(c)$. Because each level uses fixed design parameters, the actual height trails the ideal height by at most $p+n\eta$ on $A_b$ and $n\eta$ on each $H_{b,r}$. Choosing $k$ large, $p>0$ small, $N=\lceil (c+1)/p\rceil$, and $\eta$ small enough forces $J_b\mapsto A_b$ in every block (including ruling out off-support $J_b$, which still denies $A_b$); then $\disc_B\ge Np-(n-N)\eta>c$.
\end{proof}

\begin{theorem}[Two-Row Ping-Pong Cascade]\label{thm:pingpong}
For every $\eps>0$, there exists a finite support-2, strictly positive instance with $D=1$ and $2N+1$ rows forcing prefix discrepancy at least $1+\W(1)-\eps$, where $\Omega=\W(1)\approx 0.567143$ solves $\Omega e^\Omega=1$. Hence $\kappa\ge 1+\Omega\approx 1.567143$.
\end{theorem}
\begin{proof}
Fix $c\in(1,1+\Omega)$ and use only two fresh rows $\{A_b,H_b\}$ per block. By Lemma~\ref{lem:draingain}, after Level~1 ($A_b\to H_b$) from $(1,0)$, row $A_b$ ends at $c-1>0$ rather than $0$, while $H_b$ reaches $\alpha_1=-\ln(c-1)$. The remaining gap on $H_b$ is $u_1:=c-\alpha_1=1+(c-1)+\ln(c-1)<1$ precisely because $(c-1)e^{c-1}<1\iff c-1<\W(1)=\Omega$. Immediately reversing direction ($H_b\to A_b\to H_b\to\cdots$) gives the passive row a free head start $\beta=c-1$ on every level $r\ge 2$, producing the exact gap recurrence
\[
 u_r := c-\alpha_r = c - \bigl((c-1)-\ln u_{r-1}\bigr) = 1+\ln u_{r-1}\qquad (r\ge 2).
\]
On $(0,1)$, the positive decrement $\delta(u):=u-(1+\ln u)$ has derivative $1-1/u<0$, so $\delta(u)$ is strictly decreasing in $u$; hence $u_{r-1}-u_r\ge \delta(u_1)>0$ uniformly until $u_r<0$ ($\alpha_r>c$) in finitely many levels. Since every level's drain and gain are fixed design constants, the slacks remain bounded by $p+n\eta$ on $A_b$ and $n\eta$ on $H_b$ across all levels without exponential amplification.
\end{proof}

Above $1+\Omega$, after Level~1 neither $X=c-1$ nor $Y=-\ln(c-h)$ alone lies within distance $1$ of $c$, yet their sum still exceeds $1$.

\begin{lemma}[Height-Amplifying Binary Split]\label{lem:split}
At budget $c>1$, let $X$ sit at height $h$ with $u:=c-h\in(0,1)$ and let $Y$ be fresh at $0$. Run Level~1 ($X\to Y$) so $(X,Y)\to(c-1,-\ln u)$. If $u<e^{-c}$, then already $Y=-\ln u>c$ and the branch dies immediately. Otherwise $e^{-c}\le u<1$, so $|Y-X|\le 1$, and appending one weight-$(1,1)$ split column $K$ on $\{X,Y\}$ with $x_X=\frac{1+Y-X}{2}\in[0,1]$ and $x_Y=1-x_X$ guarantees that under every integral assignment (to $X$, to $Y$, or off-support), at least one of $\{X,Y\}$ reaches
\[
 h' = \frac{1+(c-1)-\ln u}{2} = \frac{c-\ln u}{2}.
\]
Attaching independent killer subtrees on $X$ and then on $Y$ (on disjoint fresh row sets) kills both branches with at most $O(n\eta)$ off-support interference, where $n$ is the full column count across all leaves.
\end{lemma}
\begin{proof}
Under $K\mapsto X$, $\disc_Y$ rises to $Y+x_Y=Y+\frac{1+X-Y}{2}=\frac{1+X+Y}{2}=h'$. Under $K\mapsto Y$, $\disc_X$ rises to $X+x_X=\frac{1+X+Y}{2}=h'$. Under an off-support assignment of $K$, both $X$ and $Y$ receive their fractional shares and reach $h'$. Because the two killer subtrees use disjoint fresh auxiliary rows, walking through the non-triggered subtree can pull the surviving sibling down only via off-support dumps totaling at most $n\eta$.
\end{proof}

\begin{theorem}[Recursive Binary-Tree Splitting + Ping-Pong]\label{thm:ln2}
For every $\eps>0$, there exists a finite support-2, strictly positive instance with $D=1$ forcing prefix discrepancy at least $1+\ln 2-\eps$. Hence $\kappa\ge 1+\ln 2\approx 1.69314718$.
\end{theorem}
\begin{proof}
Fix $c\in(1,1+\ln 2)$. In gap coordinates $u=c-h\in(0,1)$, Lemma~\ref{lem:split} maps $u\mapsto u'=(c+\ln u)/2$, so
\[
 u-u' = \frac{2u-\ln u-c}{2} \;\ge\; \frac{1+\ln 2-c}{2} \;>\; 0
\]
uniformly on $(0,1)$, because $\phi(u)=2u-\ln u$ achieves its unique global minimum $1+\ln 2$ at $u=1/2$. Starting from $u_0=c-1<1$, after finitely many split depths $L_{\mathrm{tree}}=L_{\mathrm{tree}}(c)$ every guaranteed child gap $u_{L_{\mathrm{tree}}}$ drops below the ping-pong threshold $e^{-(c-1)}$ (or dies earlier if $u_r<e^{-c}$), where a two-row ping-pong leaf killer (Theorem~\ref{thm:pingpong}) finishes each of the $2^{L_{\mathrm{tree}}}$ leaves. Down the tree, the left child inherits slack $\le p+n\eta$ and the right child only $\le n\eta$, where $\eta$ is scaled against the exact full column count $n$ of all $2^{L_{\mathrm{tree}}}$ leaves.
\end{proof}

Exact rational evaluations of the proof recursions in \texttt{detector-chains/frontier\_next/code/verify\_cascade\_and\_pingpong.py} reach $1.363999$ (fresh cascade), $1.561479$ (ping-pong), and $1.681499$ (depth-26 binary tree; note that the script evaluates the exact rational proof recursion with full $n\eta$ accounting rather than materializing the $n\approx 6.15\times 10^{14}$ tree). Independent HiGHS MILP checks on materialized single killer blocks (with $J\to A$ forbidden and off-support allowed) confirm min-max $1.285>1+\W(1/e)$ (12 cols), $1.390>1+1/e$ (29 cols), and infeasibility at $1.569>1+\W(1)$ (5 rows, 959 cols).

\section{Sharp Machine Threshold: $m^*=3$ and the Even-$m$ Ladder}

Previously, $\kappa_2=1$ (\texttt{chairman-small-rows}) and $\kappa_{11}>1$ (\texttt{chairman-counterexample}) left $m^*\in[3,11]$. Two mechanisms close this gap completely to $m^*=3$.

\begin{lemma}[Dirty-Helper Ping-Pong]\label{lem:dirty}
At budget $c\ge 1$, suppose row $A$ enters a ping-pong block on $\{A,H\}$ at height $\alpha_0$ ($g_0:=c-\alpha_0\in(0,1)$), while $H$ enters at \emph{any} legal height $\disc_H\ge -c$ (arbitrarily polluted). After Level~1 ($A\to H$), $A$ stands at $c-1$ and $H$ rises to $\ge -c-\ln g_0$. Ping-pong on $\{A,H\}$ explodes past $c$ whenever $c-(-c-\ln g_0)<1$, i.e., $g_0<e^{1-2c}$.
\end{lemma}
\begin{proof}
By Lemma~\ref{lem:draingain}, Level~1 ($A\to H$) drains $A$ to $c-1$ and raises $H$ from $\disc_H\ge -c$ to $\alpha_1\ge -c-\ln g_0$ (as $k\to\infty$). The resulting gap on $H$ is $u_1:=c-\alpha_1\le 2c+\ln g_0$, which is strictly less than $1$ precisely when $2c+\ln g_0<1\iff g_0<e^{1-2c}$. From Level~2 onward the passive row already starts at $c-1$, so the gap obeys the Theorem~\ref{thm:pingpong} recurrence $u_r=1+\ln u_{r-1}$ and crosses $0$ in finitely many levels.
\end{proof}

\begin{theorem}[Exact Threshold $m^*=3$ and Quantitative Bound on Three Machines]\label{thm:m3}
Let $c_3^*\approx 1.10172066$ be the unique root of $2c-1=e^{-(c-1)}+e^{1-2c}$. Then $\kappa_3\ge c_3^*>1$. Combined with $\kappa_2=1$, the minimal number of machines on which the Liu--Reis bound fails is $m^*=3$.
\end{theorem}
\begin{proof}
Work on three rows $\{B,A,H\}$ at budget $c\in(1,c_3^*)$.
\begin{enumerate}
\item \textbf{Zero-pollution initial split $J_0$:} place $J_0$ on $\{B,A\}$ with weights $(1,\eta)$ and splits $(u,1-u)$ where $u=1-\delta_0\in(0,1)$ ($0<\delta_0\ll 1$). Under all three possible assignments ($J_0\to B$, $J_0\to A$, $J_0\to H$), rows $A$ and $H$ move by at most $\eta$, leaving $\{A,H\}$ essentially fresh. Meanwhile, under $J_0\not\to B$, $\disc_B=+u=1-\delta_0$; under $J_0\to B$, $\disc_B=-(1-u)=-\delta_0$.
\item \textbf{Block~1 (clean ping-pong on $(B,H,A)$ with share $p_1$):} place $J_1$ on $\{B,H\}$ with weights $(1,1)$ and splits $(p_1,1-p_1)$, followed by clean ping-pong on $\{H,A\}$. Whenever $c-(1-p_1)<e^{-(c-1)}$ (i.e., $p_1<1-c+e^{-(c-1)}$), Block~1 forces $J_1\mapsto H$ on every surviving schedule. Choosing $u+p_1>c$ (i.e., $p_1>c-1+\delta_0$) kills every $J_0\not\to B$ branch immediately ($\disc_B=u+p_1>c$). The surviving $J_0\to B$ branch exits with $\disc_B=u+p_1-1=p_1-\delta_0$ and $\disc_A\ge -c$ (any legal schedule stays above $-c$).
\item \textbf{Block~2 (dirty-helper ping-pong directly on $\{B,A\}$):} with no coupling column $J_2$, run dirty-helper ping-pong directly from active row $B$ onto dirty helper $A$. The entry gap on $B$ is $g_0=c-(u+p_1-1)=c-p_1+\delta_0$. By Lemma~\ref{lem:dirty}, Block~2 kills the surviving branch whenever $c-p_1<e^{1-2c}$, i.e., $p_1>c-e^{1-2c}$ (which automatically implies $p_1>c-1$ since $e^{1-2c}<1$ for $c>1$).
\end{enumerate}
A valid $p_1$ satisfying $c-e^{1-2c}<p_1<1-c+e^{-(c-1)}$ exists if and only if $2c-1<e^{-(c-1)}+e^{1-2c}$, i.e., $c<c_3^*\approx 1.10172066$. Choosing $\delta_0,k,\eta$ sufficiently small absorbs all $O(n\eta)$ sinks and yields finite exact rational instances forcing discrepancy $>c$.
\end{proof}

\begin{remark}[Ultra-compact cleaned-helper $3\times 33$ counterexample]
At budget $c=1$, replacing Block~1 by a single 1-level chain ($k_1=10$) and appending one modulo-$1$ helper cleanup column $Z_1$ on $\{H,A\}$ with weights $(\eta,1)$ and split $x_H=\delta+1/200$ (where $\delta=\mathrm{term}_A\bmod 1$) forbids $Z_1\to A$ on the lower modulo-$1$ exit state (which would push $A<-1$) and lifts $A$ to $\ge -1/200$ on every surviving schedule; a 7-level clean ping-pong Block~2 ($k=4$) then kills the surviving branch. This yields an explicit \textbf{$3\times 33$ instance} certified in exact \texttt{fractions.Fraction} arithmetic to satisfy $\kappa_3>1001/1000$ (\texttt{detector-chains/frontier\_v2/code/verify\_m3\_and\_small\_m\_ladder.py}), confirmed independently by exhaustive exact-rational BFS over all $3^{33}$ schedules (off-support included) and by HiGHS MILP infeasibility at budget $1.00$ (\texttt{verify\_independent\_milp\_v2.py}).
\end{remark}

\begin{theorem}[Even-$m$ Zero-Pollution Ladder]\label{thm:ladder}
For every even $m=2q+2\ge 4$, let $c_m^*$ be the unique root of $e^{-(c-1)}-(c-1)=2/m$. Then $\kappa_m\ge c_m^*$. In particular,
\[
 \kappa_4\ge 1.266249,\qquad \kappa_6\ge 1.362560,\qquad \kappa_8\ge 1.412195 > 1+\frac{1}{e},
\]
and $c_m^*\uparrow 1+\W(1)\approx 1.567143$ as $m\to\infty$.
\end{theorem}
\begin{proof}
On rows $\{B,R_0,A_1,H_1,\dots,A_q,H_q\}$, place zero-pollution $J_0$ on $\{B,R_0\}$ with weights $(1,\eta)$ and split $u=1-\delta_0\in(0,1)$. Run $q$ clean ping-pong blocks on $(B,A_r,H_r)$, each transferring $p_1<1-c+e^{-(c-1)}$. Choosing $u+qp_1>c$ (i.e., $qp_1>c-1+\delta_0$) kills every $J_0\not\to B$ branch. In the surviving $J_0\to B$ branch, row $R_0$ remains completely fresh ($\ge -n\eta$) and $\disc_B=u+qp_1-1=qp_1-\delta_0$, so clean final ping-pong on $\{B,R_0\}$ kills whenever $c-(qp_1-\delta_0)<e^{-(c-1)}$, i.e., $qp_1>c-e^{-(c-1)}+\delta_0$ (which automatically exceeds $c-1$). Letting $s:=e^{-(c-1)}-(c-1)=1-c+e^{-(c-1)}$, a valid $p_1$ satisfying $c-e^{-(c-1)}<qp_1<qs$ exists iff
\[
 qs > c - e^{-(c-1)} = 1 - s \iff (q+1)s > 1 \iff e^{-(c-1)}-(c-1) > \frac{1}{q+1} = \frac{2}{m}. \qedhere
\]
\end{proof}

\section{Exact Completeness Barrier at $1+\ln 2$}

Why did recursive binary-tree splitting stall at $1+\ln 2$, right where $2u-\ln u=c$ meets the ping-pong threshold $u=e^{-(c-1)}$ tangentially at $(u,c)=(1/2,1+\ln 2)$? We prove that this wall is exact across all forced chains (even with non-monotone weights) and all $q$-way disjoint-subtree splits.

\begin{theorem}[Universal Non-Monotone Drain--Gain Optimality]\label{thm:nonmonotone}
Fix $c>1$, entry height $\alpha\in(c-1,c)$ ($u:=c-\alpha$), and any $k\ge 2$, arbitrary weights $d_2,\dots,d_k\in(0,1]$ on $X$ and $d_{Y,2},\dots,d_{Y,k}\in(0,1]$ on $Y$, and splits $x_2,\dots,x_k\in[0,1]$ such that every first deviation on $X$ reaches at least $c$. Letting $w:=c-\alpha_{\mathrm{end}}$ where $\alpha_{\mathrm{end}}=\alpha-\sum_{j=2}^k d_j(1-x_j)$, we have $w\in[u,1]$ and
\[
 \disc_Y^{\mathrm{gain}} := \sum_{j=2}^k d_{Y,j}(1-x_j) \;\le\; \ln\!\left(\frac{w}{u}\right) \;\le\; -\ln u,
 \qquad \alpha_{\mathrm{end}}+\disc_Y^{\mathrm{gain}} \;\le\; c-1-\ln u.
\]
\end{theorem}
\begin{proof}
Let $u_1:=u$ and $u_j:=u_{j-1}+d_j(1-x_j)$ for $j=2,\dots,k$, so $u_k=w$. The deviation condition at column $j$ states $(c-u_{j-1})+d_j x_j\ge c$, i.e., $d_j x_j\ge u_{j-1}$, hence $u_j\le d_j\le 1$ (giving $w\le 1$) and
\[
 d_{Y,j}(1-x_j)\le 1-x_j = \frac{d_j(1-x_j)}{d_j} \le \frac{u_j - u_{j-1}}{u_j} \le \int_{u_{j-1}}^{u_j}\frac{dt}{t} = \ln\!\left(\frac{u_j}{u_{j-1}}\right).
\]
Summing over $j=2,\dots,k$ gives $\disc_Y^{\mathrm{gain}}\le \ln(w/u)\le -\ln u$. Finally, $\alpha_{\mathrm{end}}+\disc_Y^{\mathrm{gain}}\le c-w+\ln(w/u)$, whose derivative $-1+1/w\ge 0$ on $(0,1]$ is maximized at $w=1$, giving $c-1-\ln u$.
\end{proof}

\begin{theorem}[Barrier for Fresh-Partner Chains and Abstract Disjoint Splits at $1+\ln 2$]\label{thm:barrier}
Fix $c\ge 1+\ln 2$, set $h^*:=c-e^{-(c-1)}$ (the single-row ping-pong kill threshold with a fresh partner), and define the $1$-Lipschitz row-excess potential $E(a):=\max(0,\, a-(h^*-1))$.
\begin{enumerate}
\item For any forced chain from active height $\alpha\in(c-1,h^*)$ into a fresh row $Y$ ($\beta\le 0$) satisfying Theorem~\ref{thm:nonmonotone}, the total excess change obeys
\[
 \Delta E := E(\alpha_{\mathrm{end}})+E(\disc_Y^{\mathrm{gain}})-E(\alpha) \;\le\; 0
\]
(and strictly $\Delta E < 2e^{-(c-1)}-1\le 0$ whenever $Y$'s excess activates).
\item For any $q$-way split column ($q\ge 2$) on rows $1,\dots,q$ with arbitrary weights $d_r\in(0,1]$ and splits $\sum_{r=1}^q x_r=1$, partitioned into pairwise disjoint child subtrees $I_1,\dots,I_q$ (with $r\in I_r$ allowed), if the pre-split total excess satisfies $S:=\sum_i E(a_i)<1$, let $L:=\{r\in\{1,\dots,q\}: a_r - d_r(1-x_r)\ge -c\}$ be the set of $-c$-legal branches. Whenever $L\ne\emptyset$, there exists a legal branch $r\in L$ whose child subtree satisfies
\[
 \sum_{i\in I_r} E\!\left(a_i^{(r)}\right) \;\le\;
 \begin{cases}
  \frac{S+1}{|L|} < 1 & \text{if } |L|\ge 2,\\
  S < 1 & \text{if } |L|=1.
 \end{cases}
\]
\end{enumerate}
Consequently, in the abstract gadget game where only forced fresh-partner chains ($\beta\le 0$), single-row fresh-partner leaf killers, and disjoint child-subtree branches are used (each branch facing only its own $I_r$), starting from $(1,0,\dots,0)$ (where $E(1)=2-c+e^{-(c-1)}<1$) never reaches $h^*$ at any budget $c\ge 1+\ln 2$.
\end{theorem}
\begin{proof}
(1) Write $u=c-\alpha\in(e^{-(c-1)},1)$ and $w=c-\alpha_{\mathrm{end}}\in[u,1]$. Since $h^*-1=c-1-e^{-(c-1)}$ and $w\le 1<1+e^{-(c-1)}$, both $\alpha$ and $\alpha_{\mathrm{end}}$ lie strictly above $h^*-1$, so
\[
 E(\alpha)=1+e^{-(c-1)}-u,\qquad E(\alpha_{\mathrm{end}})=1+e^{-(c-1)}-w,\qquad E(\alpha_{\mathrm{end}})-E(\alpha)=u-w\le 0.
\]
If $\disc_Y^{\mathrm{gain}}\le h^*-1$, then $E(\disc_Y^{\mathrm{gain}})=0$ and $\Delta E=u-w\le 0$. Otherwise $E(\disc_Y^{\mathrm{gain}})=\disc_Y^{\mathrm{gain}}-(h^*-1)$, and Theorem~\ref{thm:nonmonotone} gives
\[
 \Delta E \;\le\; u - w + \ln\!\left(\frac{w}{u}\right) - (h^*-1)
 \;\le\; (u - 1 - \ln u) - (c - 1 - e^{-(c-1)})
 \;=\; u - \ln u - c + e^{-(c-1)},
\]
using that $w\mapsto -w+\ln w$ is nondecreasing on $(0,1]$ and maximized at $w=1$. On $(e^{-(c-1)},1)$, $u\mapsto u-\ln u$ has derivative $1-1/u<0$ and is strictly bounded by its value at $u=e^{-(c-1)}$, yielding $\Delta E < 2e^{-(c-1)}-1\le 0$ for all $c\ge 1+\ln 2$.

(2) For any blocked row $b\notin L$, $a_b - d_b(1-x_b) < -c \implies a_b + d_b x_b < d_b - c \le 1 - c < 0 \le h^*-1$, so $E(a_b^{(r)})=0$ on every branch $r\in L$. For $r\in L$, row $r$ decreases ($E(a_r^{(r)})\le E(a_r)$), while each unblocked $i\in L\setminus\{r\}$ increases by at most $x_i$ ($E(a_i^{(r)})\le E(a_i)+x_i$). Summing over $r\in L$ on pairwise disjoint $I_r$ gives $\sum_{r\in L}\sum_{i\in I_r} E(a_i^{(r)})\le S + \sum_{r\in L}\sum_{i\in (I_r\cap L)\setminus\{r\}} x_i$, which is $\le S$ if $|L|=1$ and $\le S+1$ if $|L|\ge 2$. Averaging over $r\in L$ yields the stated bound.
\end{proof}

\begin{remark}[Exact scope of Theorem~\ref{thm:barrier}]
Theorem~\ref{thm:barrier} is a barrier for fresh-partner chains ($\beta\le 0$) and the abstract disjoint-subtree game, \emph{not} an unconditional upper bound on all offline constructions. Indeed, at $c=1.6932>1+\ln 2$, a pair of rows pre-loaded at $(0.95,0.43)$ (total excess $S=0.9935<1$, both below $h^*\approx 1.193$)---or $(1.15,0.19)$ at $c=1+\ln 2$---is immediately killed by plain two-row ping-pong because the partner row is non-fresh ($\beta>0$).
\end{remark}

\section{BFS-Box Trapping ($\kappa_m\le m-1$), Split Optimality, and Open Problems}

The standard $0/1$-simplex iterated prefix basis-rounding specialization of B\'ar\'any--Grinberg \cite{BaranyGrinberg1981}, Lenstra--Shmoys--Tardos \cite{LST1990}, and Beck--Fiala \cite{BeckFiala1981} maintains a basic feasible solution on at most $m$ fractional columns and yields $|\disc_i(t)|<mD$, i.e., $\kappa_m\le m$. We improve this baseline by a full unit for every $m\ge 1$ by coupling oriented BFS pivots with a box-trapping invariant.

\begin{theorem}[BFS-Box Trapping: $\kappa_m\le m-1$]\label{thm:km1}
For every $m\ge 1$, every $n\ge 1$, every fractional assignment $x\in[0,1]^{m\times n}$ ($\sum_{i=1}^m x_{ij}=1$), and every weight matrix $d\in[0,D]^{m\times n}$, there exists a support-preserving integral assignment $y\in\{0,1\}^{m\times n}$ ($\sum_{i=1}^m y_{ij}=1$, $y_{ij}=0$ whenever $x_{ij}=0$) satisfying
\[
 |\disc_i(t)|\le (m-1)D\qquad\text{for all }i\in\{1,\dots,m\}\text{ and }t\in\{1,\dots,n\}.
\]
Hence $\kappa_m\le m-1$ for every $m\ge 1$ (unifying $\kappa_2=1$, and tightening $1.10172\le \kappa_3\le 2$ and $1.26625\le \kappa_4\le 3$). Moreover, one such assignment is computable in $O(n\cdot\mathrm{poly}(m))$ exact arithmetic operations.
\end{theorem}
\begin{proof}
Normalize $D=1$ and set $Q:=[-(m-1),m-1]^m$. For $m=1$ every column is integral and $\disc_1(t)\equiv 0$, so assume $m\ge 2$. We maintain by induction on $t=0,1,\dots,n$ a subset $F_t\subseteq\{1,\dots,t\}$ of fractional columns with vectors $z_j^{(t)}\in\Delta_{m-1}$ ($|\mathrm{supp}(z_j^{(t)})|\ge 2$, $\mathrm{supp}(z_j^{(t)})\subseteq\mathrm{supp}(x_j)$) such that:
\begin{enumerate}
\item[(i)] \textbf{Basic Feasible Solution:} On $E_t:=\{(i,j):j\in F_t,\; z_{ij}^{(t)}>0\}$, the homogeneous system $\sum_{j:(i,j)\in E_t} d_{ij}\delta_{ij}=0$ ($1\le i\le m$), $\sum_{i:(i,j)\in E_t}\delta_{ij}=0$ ($j\in F_t$) has trivial kernel ($\delta\equiv 0$).
\item[(ii)] \textbf{In-Box Trapping:} For every integral completion $y_{F_t}\in\prod_{j\in F_t}\mathrm{supp}(z_j^{(t)})$ whose affine discrepancy $\disc^{(t)}(y_{F_t}):=\sum_{j\in F_t} d_{\cdot,j}\circ(z_j^{(t)}-y_j)$ lies in $Q$, there exists a support-preserving schedule $\sigma^{(t)}(y_{F_t})$ on columns $1,\dots,t$ whose prefixes all stay in $Q$ and whose time-$t$ discrepancy equals $\disc^{(t)}(y_{F_t})$.
\end{enumerate}

\emph{Base case ($t=0$):} $F_0=\emptyset$, $\disc^{(0)}()=0\in Q$.

\emph{Structural properties of any BFS $z^{(t)}$:} Let $R:=\{i:\deg(i)\ge 1\}$ in the bipartite support graph $(R,F_t,E_t)$. Since trivial kernel forces $|E_t|\le |R|+|F_t|$ while $\deg(j)\ge 2$ for every $j\in F_t$, we have $2|F_t|\le |E_t|\le |R|+|F_t|\le m+|F_t|$, so $|F_t|\le |R|\le m$, with equality $|F_t|=m$ forcing $|R|=m$, $|E_t|=2m$, and $\deg(j)=2$ for every $j\in F_t$. For any row $i$ and any completion $y_{F_t}$, each nonzero term has $|d_{ij}(z_{ij}^{(t)}-y_{ij})|<1$, so $|\disc_i^{(t)}(y_{F_t})|<\deg(i)$. Thus $|\disc_i^{(t)}(y_{F_t})|>m-1$ is possible \emph{only} if $\deg(i)=m$ (hence $|F_t|=m$ and every $j\in F_t$ has support $\{i,r_j\}$ of size $2$). For $m\ge 3$, at most one hub row $i^*$ can have degree $m$ (two degree-$m$ rows would exhaust all $2m$ edges and force $|R|=2<m$), and every $r\ne i^*$ has $\deg(r)\le 2\le m-1$, hence $|\disc_r^{(t)}|<m-1$.

We claim \emph{at most one completion $y_{F_t}^*$ can lie outside $Q$}. Indeed, fix a violating row $i^*$ (unique for $m\ge 3$). If $k\in\{1,\dots,m-1\}$ columns of $F_t$ assign to $i^*$ and $m-k$ away, then $-(m-1)\le -k<\disc_{i^*}^{(t)}<m-k\le m-1$, so only the two extreme completions $\mathbf{0}_{i^*}$ ($y_{i^*j}=0\;\forall j$) and $\mathbf{1}_{i^*}$ ($y_{i^*j}=1\;\forall j$) can violate $Q$. For $m\ge 3$, both cannot violate $Q$ because $\sum_{j\in F_t} z_{i^*j}^{(t)}>m-1$ and $\sum_{j\in F_t}(1-z_{i^*j}^{(t)})>m-1$ imply $m>2(m-1)\iff m<2$. For $m=2$, both rows share the same two extreme completions $(0,0)$ and $(1,1)$ while $z_{2j}^{(t)}=1-z_{1j}^{(t)}$: if $(0,0)$ violates $Q$ on either row then $\sum d_{1j}z_{1j}^{(t)}>1$ or $\sum d_{2j}(1-z_{2j}^{(t)})>1$, forcing $\sum z_{1j}^{(t)}>1$, whereas if $(1,1)$ violates $Q$ on either row then $\sum(1-z_{1j}^{(t)})>1$, contradicting $\sum z_{1j}^{(t)}+\sum(1-z_{1j}^{(t)})=2$.

\emph{Induction step ($t\to t+1$):} If $|\mathrm{supp}(x_{t+1})|=1$ ($x_{t+1}=e_{r_0}$), append $r_0$ to all schedules and keep $F_{t+1}=F_t$. Otherwise form $F':=F_t\cup\{t+1\}$ and $z^{(0)}:=(z^{(t)},x_{t+1})$. Call a completion $(y_{F_t},y_{t+1})\in\prod_{j\in F'}\mathrm{supp}(z_j^{(0)})$ \emph{bad} if $\disc^{(t+1)}(y_{F_t},y_{t+1})\in Q$ while $\disc^{(t)}(y_{F_t})\notin Q$. If no outside-$Q$ completion $y_{F_t}^*$ exists, or if its violating row $i^*$ satisfies $i^*\notin\mathrm{supp}(x_{t+1})$, there are no bad completions; pivot $z^{(0)}$ along arbitrary nonzero kernel vectors to a BFS $z^{(t+1)}$. Otherwise $0<x_{i^*,t+1}<1$:
\begin{itemize}
\item \textbf{Case (+):} $\disc_{i^*}^{(t)}(y_{F_t}^*)=\sum_{j\in F_t} d_{i^*j}z_{i^*j}^{(t)}>m-1$ ($y_{i^*j}^*=0\;\forall j\in F_t$). Any completion with $y_{i^*,t+1}=0$ keeps $\disc_{i^*}^{(t+1)}>m-1$, so the unique bad completion is $(y_{F_t}^*,e_{i^*})$. Pivot iteratively along nonzero kernel vectors $\delta$ oriented so that $\delta_{i^*,t+1}\le 0$, taking maximal step $\tau>0$ keeping $z^{(s+1)}=z^{(s)}+\tau\delta\ge 0$. At every step $s$, row $i^*$'s conservation law gives
\[
 \sum_{j\in F_t} d_{i^*j}z_{i^*j}^{(t)} = \sum_{j\in F_t} d_{i^*j}z_{i^*j}^{(s)} + d_{i^*,t+1}\underbrace{(z_{i^*,t+1}^{(s)}-x_{i^*,t+1})}_{\le 0}.
\]
If any $j_0\in F_t$ ever hit $z_{i^*j_0}^{(s)}=0$, the right-hand side would have at most $m-1$ nonzero terms $\le 1$, contradicting $>m-1$. Hence $z_{i^*j}^{(s)}>0$ for all $j\in F_t$ and $z_{i^*,t+1}^{(s)}\le x_{i^*,t+1}<1$ throughout. At the terminal BFS $z^{(t+1)}$, $|F_{t+1}|\le m<|F'|$ forces at least one column in $F'$ integral: either $z_{i^*,t+1}^{(t+1)}=0$ (since $<1$) or $z_{i^*j}^{(t+1)}=1$ for some $j\in F_t$ (since $>0$), eliminating $(y_{F_t}^*,e_{i^*})$ from $\prod_{j\in F'}\mathrm{supp}(z_j^{(t+1)})$.
\item \textbf{Case (-):} $\disc_{i^*}^{(t)}(y_{F_t}^*)=-\sum_{j\in F_t} d_{i^*j}(1-z_{i^*j}^{(t)})<-(m-1)$ ($y_{i^*j}^*=1\;\forall j\in F_t$). Every bad completion has $y_{i^*,t+1}=0$. Orienting every pivot so $\delta_{i^*,t+1}\ge 0$ maintains $z_{i^*,t+1}^{(s)}\ge x_{i^*,t+1}>0$ and
\[
 \sum_{j\in F_t} d_{i^*j}(1-z_{i^*j}^{(t)}) = \sum_{j\in F_t} d_{i^*j}(1-z_{i^*j}^{(s)}) - d_{i^*,t+1}\underbrace{(z_{i^*,t+1}^{(s)}-x_{i^*,t+1})}_{\ge 0},
\]
so no $j\in F_t$ can hit $z_{i^*j}^{(s)}=1$. At the terminal BFS $z^{(t+1)}$, either $z_{i^*,t+1}^{(t+1)}=1$ or $z_{i^*j}^{(t+1)}=0$ for some $j\in F_t$, eliminating every bad completion.
\end{itemize}
Setting $F_{t+1}:=\{j\in F':|\mathrm{supp}(z_j^{(t+1)})|\ge 2\}$ and fixing $\hat{y}_j:=z_j^{(t+1)}$ for $j\in F'\setminus F_{t+1}$, every in-box completion $y_{F_{t+1}}$ extends to a non-bad completion $(y_{F_t},y_{t+1})$ whose prefix $\disc^{(t)}(y_{F_t})$ lies in $Q$, closing the induction. At $t=n$, at least $2^{|F_n|}-1\ge 1$ completion lies in $Q$. For runtime, checking whether an outside-$Q$ corner exists at step $t$ inspects only the two extreme corners of the (at most one) degree-$m$ row in $O(m)$ time, and once the forward pass records the fixed values $\hat{y}_{F'\setminus F_{t+1}}$ at each step, picking any in-box completion of $F_n$ and unwinding backward from $t=n$ to $1$ recovers a valid schedule in $O(n\cdot\mathrm{poly}(m))$ exact arithmetic operations.
\end{proof}

\begin{corollary}[Origin-Enclosing Hull and Sharpness of BFS-Box Trapping]\label{cor:hull}
For every $m\ge 2$ and every finite instance, the set $R_t$ of reachable prefix discrepancies inside $Q=[-(m-1),m-1]^m$ satisfies $0\in\mathrm{conv}(R_t)$ for every $t=0,\dots,n$ (restoring at $c=2$ the convex-hull property that fails at $c=1$ on $m=3$). Conversely, for every $m\ge 3$ and every $c\in[1,m-1)$ (and $m=2$, $c\in[1/2,1)$), the BFS-box trapping invariant fails already at $t=m$: taking $m-1$ support-$2$ columns on $\{0,j\}$ ($j=1,\dots,m-1$) with $x_0=1-\eps$, $d=(1,1)$, followed by an $m$-th column on $\{0,1\}$ with $x_0=1-\eps$ and $d=(1,1/2)$ ($\eps=\frac{m-1-c}{m-1/2}$ for $m\ge 3$, $\eps=\frac{4}{5}(1-c)$ for $m=2$) has trivial kernel at both $t=m-1$ and $t=m$ yet produces an un-eliminable bad completion $(0,\dots,0,1)\in[-c,c]^m$.
\end{corollary}
\begin{proof}
In the BFS $z^{(t)}$ of Theorem~\ref{thm:km1}, if no completion leaves $Q$, the product measure $\prod_j z_{r_j,j}^{(t)}$ averages $\disc^{(t)}$ to $0$. If the unique outside-$Q$ completion is $\mathbf{0}_{i^*}$ (Case (+)), then $p:=(z_{i^*j}^{(t)})_{j\in F_t}\in(0,1)^m$ satisfies $\sum_{j=1}^m p_j>m-1\ge 1$, so $p\in\mathrm{conv}(\{0,1\}^m\setminus\{\mathbf{0}\})$ (whose unique separating facet is $\sum\xi_j\ge 1$), and symmetrically $p\in\mathrm{conv}(\{0,1\}^m\setminus\{\mathbf{1}\})$ in Case (-); since $\disc^{(t)}$ maps $p\mapsto 0$ affinely, $0\in\mathrm{conv}(R_t)$. The sharpness obstruction is verified directly in exact rational arithmetic in \texttt{detector-chains/frontier\_v3/code/verify\_upper\_bound\_km\_minus\_1.py}.
\end{proof}

\begin{theorem}[Balanced Split Optimality --- First Half of Detector-Chains OP3]\label{thm:splitopt}
For every $k\ge 3$ and any nondecreasing weights $0<d_2\le\dots\le d_{k-1}\le d_k=1$, letting $R=\sum_{j=2}^{k-1}d_j/d_{j+1}$ and $\tau_{\mathrm{bal}}(d)=\frac{d_2(k-1-R)+1}{1+d_2}$, the $p=0$ chain heights $\mathrm{dev}_j(x)=1-\sum_{i=2}^{j-1}d_i(1-x_i)+d_jx_j$ ($2\le j\le k$, $\mathrm{dev}_k=\mathrm{term}_A$) and $\mathrm{term}_H(x)=\sum_{i=2}^k(1-x_i)$ satisfy
\[
 \max_{x\in[0,1]^{k-1}}\min\{\mathrm{dev}_2(x),\dots,\mathrm{dev}_{k-1}(x),\mathrm{term}_A(x),\mathrm{term}_H(x)\}
 = \min\{\tau_{\mathrm{bal}}(d),\,1+d_2\}.
\]
\end{theorem}
\begin{proof}
Set nonnegative dual weights $\alpha_j=\frac{1}{d_j}-\frac{1}{d_{j+1}}\ge 0$ ($2\le j\le k-1$), $\alpha_k=1$, and $\beta=1$. Then $\sum_{j=i}^k\alpha_j=1/d_i$ cancels every $x_i$ identically, yielding
\[
 \sum_{j=2}^{k-1}\!\left(\frac{1}{d_j}-\frac{1}{d_{j+1}}\right)\mathrm{dev}_j(x)+\mathrm{term}_A(x)+\mathrm{term}_H(x)
 \equiv \frac{1}{d_2}+k-1-R.
\]
Dividing by $\frac{1}{d_2}+1$ gives $\min\le\tau_{\mathrm{bal}}(d)$, while $\mathrm{dev}_2(x)=1+d_2 x_2\le 1+d_2$. Matching primal attainment uses $x_{j+1}=d_j/d_{j+1}$ and $x_2=\min(1,(\tau_{\mathrm{bal}}(d)-1)/d_2)$. (Optimality of chains among \emph{all} single-block gadgets---the second half of OP3---remains open.)
\end{proof}

\subsection*{Open Problems}
\begin{enumerate}
\item \textbf{Is $\kappa$ finite?} Now squeezed to $1+\ln 2\le\kappa\le\infty$. Breaking $1+\ln 2$ requires going beyond fresh-partner chains and abstract disjoint-subtree splits (e.g., pre-loaded non-fresh partners or genuinely non-disjoint multi-row coupling; Theorem~\ref{thm:barrier}).
\item \textbf{Exact value of $\kappa_3$} (currently squeezed to $1.10172\le \kappa_3\le 2$ by Theorems~\ref{thm:m3} and~\ref{thm:km1}) and asymptotic growth of $\kappa_m$ as $m\to\infty$ (beating $\kappa_m\le m-1$ must go beyond BFS-box trapping by Corollary~\ref{cor:hull}).
\item \textbf{Morell--Skutella (2022) costless two-sided unsplittable flow:} still open (unrelated to machine-dependent weights, which assign edge-by-commodity weights).
\end{enumerate}

\section*{Reproducibility and Disclosure}
All lower-bound constructions, BFS-box trapping checks, and split-optimality proofs carry exact \texttt{fractions.Fraction} verifiers in \texttt{detector-chains/frontier\_next/code/}, \texttt{detector-chains/frontier\_v2/code/}, and \texttt{detector-chains/frontier\_v3/code/}, complemented by independent HiGHS MILP infeasibility checks (allowing off-support assignments) and numerical barrier-audit scripts. AI systems (Gemini 4 Neon, Claude Opus 5.5, Claude Fable 5.2, GPT-5.6 Pro) assisted with proof discovery, verification code, and adversarial cross-model refereeing; the author assumes responsibility for independent verification and final text.

\begin{thebibliography}{10}
\bibitem{Tijdeman1973} R.~Tijdeman, On a distribution problem in finite and countable sets, \emph{J.\ Combin.\ Theory Ser.\ A} 15 (1973), 129--137.
\bibitem{Meijer1973} H.~G.~Meijer, On a distribution problem in finite sets, \emph{Indag.\ Math.} 35 (1973), 9--17.
\bibitem{Tijdeman1980} R.~Tijdeman, The chairman assignment problem, \emph{Discrete Math.} 32 (1980), 323--330.
\bibitem{BeckFiala1981} J.~Beck and T.~Fiala, ``Integer-making'' theorems, \emph{Discrete Appl.\ Math.} 3 (1981), 1--8.
\bibitem{BaranyGrinberg1981} I.~B\'ar\'any and V.~S.~Grinberg, On some combinatorial questions in finite-dimensional spaces, \emph{Linear Algebra Appl.} 41 (1981), 1--9.
\bibitem{LST1990} J.~K.~Lenstra, D.~B.~Shmoys, and \'E.~Tardos, Approximation algorithms for scheduling unrelated parallel machines, \emph{Math.\ Program.} 46 (1990), 259--271.
\bibitem{LundMagazinov2015} B.~Lund and A.~Magazinov, On a generalization of the chairman assignment problem, arXiv:1510.04536 (2015).
\bibitem{MorellSkutella2022} S.~Morell and M.~Skutella, Single source unsplittable flows with arc-wise lower and upper bounds, \emph{Math.\ Program.} (2022).
\bibitem{LiuReis2026} S.~Liu and V.~Reis, Weighted chairman assignment and flow-time scheduling, \emph{ITCS 2026} (arXiv:2511.18546).
\end{thebibliography}

\end{document}
